% =====================================================================
% Chapter 5 -- Preprocessing
% =====================================================================
\documentclass[../preamble.tex]{subfiles}
\begin{document}

\chapter{Preprocessing}
\label{ch:preprocessing}

\section{Overview}
\label{sec:preproc-overview}

The preprocessing pipeline is fixed except for one variable: the denoising
filter. The pipeline receives a raw image, applies a contrast-enhancement
step, applies either SRAD or Gaussian denoising, and normalises the result for
the chosen backbone. A controlled comparison of the two denoisers is what
makes it possible to attribute any performance difference to that single
choice.

\begin{figure}[htbp]
  \centering
  \begin{tikzpicture}[
      node distance=0.7cm,
      step/.style={draw, rectangle, minimum height=0.9cm, minimum width=2.4cm, align=center, fill=blue!8},
      arr/.style={-{Latex[length=2mm]}},
    ]
    \node[step] (raw) {Raw image};
    \node[step, right=of raw] (resize) {Bicubic resize\\(224\,\texttimes\,224)};
    \node[step, right=of resize] (clahe) {CLAHE\\\textit{fixed}};
    \node[right=of clahe] (fork) {fork};
    \node[step, below right=of fork, xshift=0.7cm] (srad) {SRAD};
    \node[step, below left=of fork, xshift=-0.7cm] (gauss) {Gaussian};
    \node[step, right=of fork, yshift=-1.4cm] (zscore) {Z-score};
    \node[step, right=of zscore] (ready) {Tensor\\$(3,224,224)$};
    \draw[arr] (raw) -- (resize);
    \draw[arr] (resize) -- (clahe);
    \draw[arr] (clahe) -- (fork);
    \draw[arr] (fork) -- (srad);
    \draw[arr] (fork) -- (gauss);
    \draw[arr] (srad) |- (zscore);
    \draw[arr] (gauss) |- (zscore);
    \draw[arr] (zscore) -- (ready);
  \end{tikzpicture}
  \caption{Preprocessing pipeline. CLAHE is fixed in both ablation arms; only
           the denoising step (SRAD vs Gaussian) is varied.}
  \label{fig:preproc-pipeline}
\end{figure}

\section{Resize and CLAHE}
\label{sec:preproc-resize-clahe}

All images are resized bicubically to $224 \times 224$ to match the
ImageNet-pretrained input of the backbone. CLAHE is applied per channel with
clip-limit $2.0$ and a tile grid of $8 \times 8$. CLAHE is treated as a
fixed step rather than as an ablation axis: the rationale is that the
ultrasound corpus is consistently low-contrast, and adding a CLAHE-on versus
CLAHE-off arm would mix contrast change with noise change.

\section{SRAD versus Gaussian}
\label{sec:preproc-srad}

Ultrasound speckle is multiplicative and signal-dependent. Perona--Malik
anisotropic diffusion uses a gradient-magnitude diffusion coefficient that
assumes additive noise and is therefore a poor fit for speckle. SRAD
introduces the instantaneous coefficient of variation:
\begin{equation}
  q^2(x,y;t) = \frac{(\nabla^2 I(x,y;t))^2}{I(x,y;t)^2 + \text{neighbour}^2 + \epsilon},
  \label{eq:icov}
\end{equation}
which is the natural statistic for multiplicative noise. The iteration
continues until a fixed iteration budget is reached; the values are
normalised to $[0,1]$ during iteration to stabilise the numerics.

Gaussian blur is included as a control. It removes similar amounts of
high-frequency content but does not adapt to the local coefficient of
variation. Any difference between the two denoisers is therefore evidence
that SRAD's content-aware behaviour matters, not just smoothing.

\section{Letterbox padding}
\label{sec:preproc-padding}

Aspect ratios vary across the source and target datasets. The default
behaviour is to bicubic-resize the image to fill the $224 \times 224$ square
without preserving aspect ratio. An alternative mode is letterbox padding,
which preserves the aspect ratio by inserting a constant-valued border.
Letterbox padding can introduce a wide dark margin that the model can use as
a shortcut. The thesis therefore uses the no-pad variant in the reported
fine-tuned runs and treats padding as an ablation variable in the
foundation-stage sweep.

\begin{figure}[htbp]
  \centering
  \includegraphics[width=0.7\linewidth]{figures/f25_padding_vs_nopad_external_auc.png}
  \caption{Letterbox padding effect on PCOSgen AUC. Removing padding improves
           external AUC by approximately $+0.25$ absolute across the
           foundation-stage models.}
  \label{fig:padding-vs-nopad}
\end{figure}

\section{Normalisation}
\label{sec:preproc-norm}

By default the pipeline applies per-channel Z-score normalisation with mean
and standard deviation computed at the per-image level. For transfer learning
with ImageNet-pretrained backbones, ImageNet statistics
$(0.485, 0.456, 0.406)$ for mean and $(0.229, 0.224, 0.225)$ for standard
deviation are also available; the executed fine-tuned runs use the default
Z-score scheme.

\section{Augmentation}
\label{sec:preproc-aug}

The source corpus is heavily compressed by a lossy JPEG codec. A quantisation
audit of the source set found a median estimated quality factor of
approximately 40. JPEG re-compression at a random quality in $[50, 95]$ is
applied with probability $0.3$ per training sample, and a Gaussian blur
$\sigma \in [0.1, 1.5]$ is applied with probability $0.2$. Both augmentations
are stochastic and applied at training time only; validation and test splits
receive un-augmented inputs.

\section{Foundation-stage preprocessing ablation}
\label{sec:preproc-foundation-ablation}

The foundation stage ran the same nine backbones twice: once with SRAD and
once with Gaussian. The 18-run sweep produced the per-prep-per-arch external
AUC matrix that establishes the foundation checkpoint behaviour. The mean
SRAD-versus-Gaussian $\Delta\mathrm{AUC}$ on the external set is $0.022$ in
favour of SRAD. This is small in absolute terms but is consistent across
architectures, and it is the basis for choosing SRAD as the default
preprocessing in the fine-tune and ensemble stages.

\begin{figure}[htbp]
  \centering
  \includegraphics[width=0.85\linewidth]{figures/f2_external_auc.png}
  \caption{External PCOSgen AUC at the foundation stage. Both denoisers
           collapse on the external test set; the fine-tune stage is
           required to recover.}
  \label{fig:external-collapse}
\end{figure}

\end{document}